\chapter{Options and Bond Arithmetic at Speed}\label{ch:iv2:options-and-bond-arithmetic-at-speed}

\begin{ivtrains}
Numbers an interviewer expects in your head: at-the-money option values and Greeks without a calculator, parity by hand, DV01 and yield moves, carry, and ETF and futures fair value. Say the estimate first, then how far off it can be. Rates are simple unless a stem says otherwise.
\end{ivtrains}

\begin{question}[$\star$ \iqroles{QT, Trader, QR} \iqfirm{market maker, hedge fund, bank} \iqlevel{junior}]\label{iq:iv2:options-and-bond-arithmetic-at-speed:1}
A stock trades at \$100 with a volatility of 20\% a year, and interest rates are zero. Without a calculator, what is a one-year at-the-money call worth? What about a six-month at-the-money call on a \$50 stock with 30\% volatility? Give the rule you use and say roughly how accurate it is.
\end{question}

\begin{question}[$\star$ \iqroles{QT, Trader} \iqfirm{market maker, hedge fund} \iqlevel{junior}]\label{iq:iv2:options-and-bond-arithmetic-at-speed:2}
A stock trades at \$200 with an implied volatility of 40\%. Rates are zero. What does a three-month at-the-money straddle, a call plus a put at the same strike, cost? How far must the stock move for it to break even at expiry? What is one typical daily move, in percent, with 252 trading days in a year?
\end{question}

\begin{question}[$\star$ \iqroles{QT, Trader, QR} \iqfirm{market maker, hedge fund} \iqlevel{junior}]\label{iq:iv2:options-and-bond-arithmetic-at-speed:3}
Rates are zero. An at-the-money one-year call on a stock with 20\% volatility has what delta, in the Black-Scholes model? And the put at the same strike? Estimate them by hand, and say why the call's delta is not exactly one half.
\end{question}

\begin{question}[$\star$ \iqroles{QT, Trader} \iqfirm{market maker, hedge fund, bank} \iqlevel{junior}]\label{iq:iv2:options-and-bond-arithmetic-at-speed:4}
A \$100 stock has at-the-money options. By how much does a one-year at-the-money call change in value when implied volatility rises one percentage point? What about a three-month one? And a one-year call on a \$500 stock? Estimate each by hand, and say what the number depends on.
\end{question}

\begin{question}[$\star$ \iqroles{QT, Trader, QR} \iqfirm{market maker, hedge fund, bank} \iqlevel{junior}]\label{iq:iv2:options-and-bond-arithmetic-at-speed:5}
A stock trades at \$100 and pays no dividends. The one-year interest rate is 5\%, compounded continuously. A one-year call struck at \$100 costs \$12. What must the one-year put at the same strike cost? Give a quick estimate and the exact value.
\end{question}

\begin{question}[$\star$ \iqroles{QT, Trader, Risk} \iqfirm{bank, hedge fund} \iqlevel{junior}]\label{iq:iv2:options-and-bond-arithmetic-at-speed:6}
You hold \$25 million of a bond with a modified duration of 6.5 and a convexity of 60. Yields rise 30 basis points, or 0.30 percentage points. What is your loss from duration alone? What does convexity change it to?
\end{question}

\begin{question}[$\star$ \iqroles{QT, Trader, Risk} \iqfirm{bank, hedge fund} \iqlevel{junior}]\label{iq:iv2:options-and-bond-arithmetic-at-speed:7}
You hold \$10 million market value of a bond with a modified duration of 8.2. What is its DV01, the change in value for a one basis point move in yield? A bond future has a DV01 of \$65 per contract. How many contracts do you sell to hedge?
\end{question}

\begin{question}[$\star$ \iqroles{QT, Trader, QR} \iqfirm{bank, hedge fund} \iqlevel{junior}]\label{iq:iv2:options-and-bond-arithmetic-at-speed:8}
A zero-coupon bond pays 100 in 10 years, and its yield is 4\%, compounded annually. What is its price, by hand? What are its Macaulay and modified durations? What is its price change, roughly, if its yield rises to 4.1\%?
\end{question}

\begin{question}[$\star$ \iqroles{QT, Trader, QR} \iqfirm{bank, hedge fund} \iqlevel{junior}]\label{iq:iv2:options-and-bond-arithmetic-at-speed:9}
The one-year zero rate is 3\%, and the two-year zero rate is 4\%, both compounded annually. What is the one-year rate, one year forward? Give the quick answer and the exact one, and say why the quick answer works.
\end{question}

\begin{question}[$\star$ \iqroles{QT, Trader} \iqfirm{bank, hedge fund, market maker} \iqlevel{junior}]\label{iq:iv2:options-and-bond-arithmetic-at-speed:10}
Spot is 1.10 dollars per euro. The one-year dollar rate is 5\%, and the one-year euro rate is 3\%, both simple. What is the one-year forward rate, quick and exact? If someone quotes a forward of 1.13, how do you lock in a profit?
\end{question}

\begin{question}[$\star$ \iqroles{QT, Trader} \iqfirm{market maker, hedge fund} \iqlevel{junior}]\label{iq:iv2:options-and-bond-arithmetic-at-speed:11}
An ETF's shares are created in units of 50\,000. The basket behind one ETF share costs \$49.90 to buy. It costs \$0.03 per ETF share, in total, to buy the basket and to create. The ETF is quoted at \$50.02 bid and \$50.04 ask. What is your trade, and what do you make per creation unit?
\end{question}

\begin{question}[$\star$ \iqroles{QT, Trader} \iqfirm{market maker, hedge fund} \iqlevel{junior}]\label{iq:iv2:options-and-bond-arithmetic-at-speed:12}
A \$100 stock's three-month at-the-money straddle trades at \$12, and rates are zero. What implied volatility does that imply, by hand? What is the one-standard-deviation move to expiry in dollars?
\end{question}

\begin{question}[$\star\star$ \iqroles{QT, Trader, QR} \iqfirm{market maker, hedge fund} \iqlevel{junior}]\label{iq:iv2:options-and-bond-arithmetic-at-speed:13}
Rates are zero. An at-the-money one-year call on a \$100 stock with 20\% volatility is worth about \$8. By hand, estimate its gamma, the change in delta per \$1 move in the stock. Estimate its one-day theta, with 252 trading days in a year, in two ways. How much must the stock move in one day for gamma's gain to pay that theta?
\end{question}

\begin{question}[$\star\star$ \iqroles{QT, Trader} \iqfirm{market maker, hedge fund, bank} \iqlevel{junior}]\label{iq:iv2:options-and-bond-arithmetic-at-speed:14}
A stock trades at \$100 and pays a \$2 dividend in three months. Interest is 4\% a year, compounded continuously. Six-month European options have a strike of \$95. What is the call price minus the put price, by hand? What would it be if you forgot the dividend?
\end{question}

\begin{question}[$\star\star$ \iqroles{QT, Trader} \iqfirm{market maker, hedge fund} \iqlevel{junior}]\label{iq:iv2:options-and-bond-arithmetic-at-speed:15}
A stock trades at \$100, and six-month European options struck at \$100 have a call price of \$5.50 and a put price of \$5.00. The six-month interest rate is 3\% a year, simple. What forward price do the options imply? What does carry alone imply? What is the gap, and what could it stand for?
\end{question}

\begin{question}[$\star\star$ \iqroles{QT, Trader, Risk} \iqfirm{bank, hedge fund} \iqlevel{junior}]\label{iq:iv2:options-and-bond-arithmetic-at-speed:16}
A 10-year bond pays a 5\% annual coupon and yields 5\%. Why is its price 100, and what is a quick formula for its modified duration? Estimate its price if its yield rises to 5.1\%, and compare with the exact price.
\end{question}

\begin{question}[$\star\star$ \iqroles{QT, Trader} \iqfirm{bank, hedge fund} \iqlevel{junior, experienced}]\label{iq:iv2:options-and-bond-arithmetic-at-speed:17}
A 5-year bond yields 4.5\% with a modified duration of 4.5. You finance it at 3.5\% in the repurchase market. The yield curve near it slopes up by 10 basis points per year of maturity, and it stays the same over the next three months. What are your carry, your roll-down and your total, in price points, over the three months? How far can its yield rise before you lose money?
\end{question}

\begin{question}[$\star\star$ \iqroles{QT, Trader} \iqfirm{market maker, hedge fund} \iqlevel{junior, experienced}]\label{iq:iv2:options-and-bond-arithmetic-at-speed:18}
An ETF holds only Japanese stocks and trades in the US while Tokyo is closed. Its published net asset value, from Tokyo's last close, is \$40.00. Since that close, futures on the same index have risen 1\%, and the currency has not moved. The ETF trades at \$40.50. What premium does the published value show? What is the real one? Your round-trip cost is 20 basis points: is it a trade?
\end{question}

\begin{question}[$\star\star$ \iqroles{QT, Trader} \iqfirm{market maker, hedge fund, bank} \iqlevel{junior}]\label{iq:iv2:options-and-bond-arithmetic-at-speed:19}
A stock index stands at 4\,000. Its three-month future expires in exactly a quarter of a year. Interest rates are 5\% a year, the dividend yield is 2\% a year, and you use simple carry. What is fair value? The future trades at 4\,036. What is your trade, what do you make in basis points of spot, and what remains if entry and exit together cost you 10 basis points?
\end{question}

\begin{question}[$\star\star$ \iqroles{QT, Trader, QR} \iqfirm{market maker, hedge fund} \iqlevel{junior, experienced}]\label{iq:iv2:options-and-bond-arithmetic-at-speed:20}
Rates are zero. A stock trades at \$100 with 20\% volatility. A one-year call spread, long the 100 call and short the 110 call, costs about \$3.67. What does that price tell you about the chance of finishing above about \$105, and why? What does the Black-Scholes model give for that chance?
\end{question}

\begin{question}[$\star\star$ \iqroles{QT, Trader, QR} \iqfirm{market maker, hedge fund} \iqlevel{junior, experienced}]\label{iq:iv2:options-and-bond-arithmetic-at-speed:21}
Rates are zero. A stock trades at \$100. A one-year call struck at \$110 has an implied volatility of 20\%. Treat the stock's move as normal, with a standard deviation of \$20, and use $\Phi(0.5) \approx 0.69$ and a normal density of about 0.35 at 0.5. What do you estimate? How does it compare with the Black-Scholes value, and why is it lower?
\end{question}

\begin{question}[$\star\star$ \iqroles{QT, QR, Trader} \iqfirm{market maker, hedge fund} \iqlevel{junior, experienced}]\label{iq:iv2:options-and-bond-arithmetic-at-speed:22}
Rates are zero. A stock trades at \$100. What delta does the rule of thumb give for a call struck one standard deviation above the stock, where the standard deviation is $\sigma\sqrt{T}$ in percent? What is the Black-Scholes delta if the option has one year and 20\% volatility, struck at 120? If it has a short life with a total volatility of 5\%, struck at 105? Why do they differ?
\end{question}

\begin{question}[$\star\star$ \iqroles{QT, Trader, Risk} \iqfirm{bank, hedge fund} \iqlevel{junior, experienced}]\label{iq:iv2:options-and-bond-arithmetic-at-speed:23}
Per \$1 million of face value, the DV01s of a 2-year, 5-year and 10-year bond are \$190, \$460 and \$830. You buy \$10 million of the 5-year and sell both others, with half of the risk on each wing and zero net DV01. How much of each do you sell? Yields are 4.0\%, 4.4\% and 4.6\%. What is the butterfly's spread, twice the belly's yield less the wings'? What do you make or lose if only the 5-year yield falls 5 basis points, and if all three rise 10?
\end{question}

\begin{question}[$\star\star\star$ \iqroles{QT, Trader, QR} \iqfirm{market maker, hedge fund} \iqlevel{junior, experienced}]\label{iq:iv2:options-and-bond-arithmetic-at-speed:24}
Rates are zero. You sell an at-the-money one-month straddle on a \$100 stock at 30\% implied volatility and you delta-hedge it often, using deltas at 30\%. The stock then moves about 1\% a day with no trend, with 252 trading days in a year. What do you collect? What volatility does a move of 1\% a day represent? What is your expected profit, and why is that not simply vega times the difference in volatility from a constant gamma?
\end{question}

\begin{question}[$\star\star\star$ \iqroles{QT, Trader} \iqfirm{market maker, hedge fund} \iqlevel{junior, experienced}]\label{iq:iv2:options-and-bond-arithmetic-at-speed:25}
Rates are zero. A \$100 stock reports earnings in five trading days. The at-the-money straddle that expires just after the report costs \$8, and on ordinary days the stock has a volatility of 30\%, with 252 trading days in a year. What total standard deviation does the straddle imply? How big a move does it imply on the earnings day alone? What would you do if you thought the earnings move would be about 6\%?
\end{question}
